During the radiation dominated era in the early Universe, the scale factor of
the Universe $a \propto t^{1/2}$, where $t$ is the time since the Big Bang.
During most of this era, neutrons (n) and protons (p) remain in thermal
equilibrium with each other via weak interactions. The number density ($N$) of
free neutrons or protons is related to the temperature $T$ and their
corresponding masses $m$ such that
\[ N \propto m^{3/2}\exp\left(-\frac{mc^2}{k_B T}\right), \]
as long as time $t \leq t_{\mathrm{wk}} = 1.70$\,s, when
$k_B T \geq k_B T_{\mathrm{wk}} = 800$\,keV. After $t_{\mathrm{wk}}$, the weak
interactions can no longer maintain such equilibrium, and free neutrons decay
to protons with a half-life time of $610.4$\,s.

\begin{description}
\item[(T10.1)] Let the number density of protons be $N_{\mathrm{p}}$, and that
  of neutrons be $N_{\mathrm{n}}$. Calculate the relative abundance of
  neutrons given the ratio
  $X_{\mathrm{n,wk}} = N_{\mathrm{n}}/(N_{\mathrm{n}} + N_{\mathrm{p}})$ at
  time $t_{\mathrm{wk}}$. \hfill[4]
\item[(T10.2)] Photons maintain thermal equilibrium and retain a blackbody
  spectrum at all epochs.
  \begin{description}
  \item[(T10.2a)] Find the index $\beta$, such that $T(a) \propto a^\beta$.
    \hfill[2]
  \item[(T10.2b)] Identify which of the following graphs shows the correct
    behaviour of the spectral energy density for two temperatures $T_1$ and
    $T_2$. Tick the correct option in the Summary Answersheet. \hfill[2]
  \end{description}

  \ioaafig{2025_TH_T10-f1.pdf}

\item[(T10.3)] After $t_{\mathrm{wk}}$, the process of formation of deuterium
  from protons and neutrons is governed by the Saha equation, given by the
  Indian physicist Prof.\ Meghnad Saha, which can be simplified to
  \[ \frac{N_{\mathrm{D}}}{N_{\mathrm{n}}}
     = 6.5\,\eta\left(\frac{k_B T}{m_{\mathrm{n}}c^2}\right)^{3/2}
       \exp\left(-\frac{(m_{\mathrm{D}} - m_{\mathrm{p}}
       - m_{\mathrm{n}})c^2}{k_B T}\right). \]
  Here, the baryon-to-photon ratio $\eta$ is $6.1 \times 10^{-10}$, and
  $N_{\mathrm{D}}$ is the number density of deuterium.
  \begin{description}
  \item[(T10.3a)] Plot the ratio $N_{\mathrm{D}}/N_{\mathrm{n}}$ on the grid
    in the Summary Answersheet, for at least 4 reasonably spaced values of
    temperature that lie in the domain $k_B T = [60, 70]$\,keV, and draw a
    smooth curve passing through these points. \hfill[5]
  \item[(T10.3b)] From the plot find $k_B T_{\mathrm{nuc}}$ (in keV) when
    $N_{\mathrm{D}} = N_{\mathrm{n}}$. \hfill[1]
  \item[(T10.3c)] Instead, now assume that all the free neutrons combine
    instantaneously with the protons at $k_B T_{\mathrm{nuc}}$ to form
    deuterium, all of which immediately gets converted to helium
    ($^4_2$He). Compute the corresponding epoch or time of nucleosynthesis,
    $t_{\mathrm{nuc}}$ (in s), for the formation of helium. \hfill[4]
  \end{description}
\item[(T10.4)] Calculate the value of $X_{\mathrm{n,nuc}}$ immediately before
  $t_{\mathrm{nuc}}$. \hfill[5]
\item[(T10.5)] The primordial helium abundance, $Y_{\mathrm{prim}}$, is
  defined to be the fraction of total mass in the Universe that is bound in
  helium at $T_{\mathrm{nuc}}$. Obtain a theoretical estimate for the value of
  $Y_{\mathrm{prim}}$. For the purpose of this calculation alone, assume
  $m_{\mathrm{p}} \approx m_{\mathrm{n}}$ and that the mass of helium,
  $m_{\mathrm{He}} \approx 4m_{\mathrm{n}}$. \hfill[3]
\item[(T10.6)] The primordial abundance of helium is very difficult to
  measure, as stars continuously convert hydrogen to helium in the Universe.
  The amount of processing by stars in a galaxy is characterised by the
  relative number density of oxygen (which is only produced by stars) to
  hydrogen, denoted as (O/H), in the galaxy. A compilation of the measurements
  of (O/H) and the helium abundance, $Y$, for different galaxies is plotted
  below. Use all of the points in this plot (which is reproduced in the
  Summary Answersheet) to answer the following.

  \ioaafig{2025_TH_T10-f2.pdf}

  \begin{description}
  \item[(T10.6a)] Estimate $Y$ for a blue compact dwarf galaxy with a value of
    $\mathrm{(O/H)} = 1.75 \times 10^{-4}$. \hfill[2]
  \item[(T10.6b)] Obtain the slope $dY/d\mathrm{(O/H)}$ of the straight line
    fit to the above data. \hfill[2]
  \item[(T10.6c)] Estimate the primordial helium abundance,
    $Y^{\mathrm{obs}}_{\mathrm{prim}}$, based on the above observations.
    \hfill[2]
  \end{description}
\item[(T10.7)] The deviation between $Y_{\mathrm{prim}}$ and
  $Y^{\mathrm{obs}}_{\mathrm{prim}}$ can be reconciled by changing the
  baryon-to-photon ratio $\eta$. When $\eta$ is decreased, as indicated by
  $\downarrow$ in the Summary Answersheet, indicate the increase ($\uparrow$)
  or decrease ($\downarrow$) in $N_{\mathrm{D}}/N_{\mathrm{n}}(T)$,
  $T_{\mathrm{nuc}}$ (when $N_{\mathrm{D}} = N_{\mathrm{n}}$),
  $t_{\mathrm{nuc}}$, $X_{\mathrm{n,nuc}}$, and $Y_{\mathrm{prim}}$ in the
  boxes provided in the Summary Answersheet. \hfill[3]
\end{description}
